% c03 —— 第 3 节末（Hamilton 方程、相空间）+ 第 4–7 节，p009–p013

Legendre 变换之前的 $(q,\dot q)$ 住在位形空间 $N$ 的切丛 $TN$ 中，而 Legendre 变换之后的 $(q,p)$ 住在余切丛 $T^*N$ 中。

运动方程可用 Hamilton 函数写成
\begin{equation}
    \dot p_i = -\frac{\partial H}{\partial q_i},\qquad
    \dot q_i = \frac{\partial H}{\partial p_i} \tag{3.5}
\end{equation}
这是一阶常微分方程，而且关于广义坐标 $q$ 与广义动量 $p$ 呈（反）对称的形式。这一方程的背后是 $T^*N$ 上的辛形式 $\omega=\sum_i dp_i\wedge dq_i$。相应的 Poisson 括号由
\begin{equation}
    \{f,g\}=\sum_i\left(\frac{\partial f}{\partial q_i}\frac{\partial g}{\partial p_i}-\frac{\partial f}{\partial p_i}\frac{\partial g}{\partial q_i}\right) \tag{3.6}
\end{equation}
给出，运动方程可改写为
\begin{equation}
    \dot p_i=\{H,p_i\},\qquad \dot q_i=\{H,q_i\} \tag{3.7}
\end{equation}
这一运动方程可象征地写作
\begin{equation}
    \frac{\dd}{\dd t}=\{H,-\}, \tag{3.8}
\end{equation}
或者更确切地说，对给出系统时间演化的微分同胚 $\phi\colon\R\to\mathrm{Diff}(T^*N)$，也可写作
\begin{equation}
    \dot\phi=X_H \tag{3.9}
\end{equation}

相空间（phase space）$T^*N$ 的维数是位形空间 $N$ 的两倍。位形空间的微分同胚诱导相空间的辛同胚，但在相空间的辛同胚中，来自位形空间微分同胚的只占极小一部分。

$T^*N$ 的任意函数都决定 $T^*N$ 的一个辛同胚，这是 $N$ 或 $TN$ 所不具有的、$T^*N$（或一般的辛流形，乃至更一般的 Poisson 流形）的显著特征（或特长）。

\section{性质与结构}

「是拓扑流形」是拓扑空间的一种性质（property）。另一方面，「可微（differentiable）」一词听上去像是性质，但「是微分流形」则是拓扑流形上的一种结构（structure）。拓扑流形定义为局部与 Euclid 空间同胚的拓扑空间。微分流形则定义为：在拓扑流形上以数据的形式附加一个图册，使坐标变换可微。\footnote{微分构造不是性质而是结构，这是 Milnor 的伟大发现。例如，7 维球面上的微分结构全体构成一个 28 阶循环群，而且它们全都可以（例如作为 Brieskorn 流形）具体构造出来，这是令人惊异的。}

复结构与辛结构不是微分流形的性质，而是结构。

\begin{exercise}
证明：复结构是结构，即一般地，微分流形容许多个复结构。
\end{exercise}

\begin{exercise}
证明：辛结构是结构，即一般地，微分流形容许多个辛结构。
\end{exercise}

给定的微分流形拥有多少个复结构或辛结构，一般是非常困难的问题。

\begin{exercise}
举出一个只容许唯一复结构的微分流形的例子。
\end{exercise}

\begin{exercise}
举出一个只容许唯一辛结构的微分流形的例子。
\end{exercise}

复结构与辛结构既可以用「具有标准结构的开集的粘贴」来定义，也可以用切丛的张量的语言来定义；而 Riemann 结构无法用标准结构的粘贴来定义。这一差别来源于可积性的有无。

复代数簇（complex algebraic variety）定义为复数域上约化、不可约、分离的有限型概形。概形是局部同构于仿射概形的环空间，而环空间是拓扑空间与其上环的层的组。\footnote{也有把概形不看作用环空间、而看作从环的范畴到集合的范畴的函子的观点，这里不涉及。}作为环空间，是拓扑空间上的一种结构；作为代数簇，则是环空间的一种性质。光滑代数簇自然地决定一个微分流形，因此也可以说，作为代数簇是微分流形上的一种结构。代数结构恐怕无法像复结构或辛结构那样，用张量场的语言来改写。

\begin{exercise}
举出一组不同构的代数簇，它们作为复流形却是同构的。
\end{exercise}

\section{G 结构的几何}

1872 年，在 Erlangen 大学的教授就任演讲中，Klein 提出如下观点：对空间与作用于其上的群这一组对象，相应地确定一门几何学，只有在该群作用下不变的空间性质才成为几何学的研究对象。这既把 Euclid 几何与仿射几何定位到以射影几何为顶点的层级之中，也促进了 Lie 群与齐性空间的研究。

Cartan 的 G 结构概念，把 Klein 的观点扩展为包含 Riemann 几何的形式。对 Euclid 空间 $\R^n$ 与作用于其上的 Lie 群 $G$ 的组 $(G,\R^n)$，$n$ 维流形 $M$ 的 G 结构（G-structure）是指 $M$ 的切丛 $TM$ 的结构群 $\mathrm{GL}(n,\R)$ 到 $G$ 的约化。也就是说，取 $M$ 的开覆盖，把切丛 $TM$ 表示为平凡丛的粘贴时，让转换函数取值于 $G$ 而不是 $\mathrm{GL}(n,\R)$。

把 $\mathrm{GL}(n,\R)$ 中包含单位元的连通分支记作 $\mathrm{GL}^+(n,\R)$。

\begin{exercise}
证明：$\mathrm{GL}(n,\R)/\mathrm{GL}^+(n,\R)\cong\Z/2\Z$。
\end{exercise}

\begin{exercise}
证明：$\mathrm{GL}^+(n,\R)$ 结构与定向等价。
\end{exercise}

\begin{exercise}
证明：$\mathrm{O}(n)$ 结构与 Riemann 度量等价。
\end{exercise}

\begin{exercise}
证明：$\mathrm{GL}(n/2,\C)$ 结构与近复结构等价。
\end{exercise}

G 结构往往由张量场来刻画；在上述三个例子里，分别成为体积形式（处处不为零的 $n$ 次微分形式）、Riemann 度量（处处正定的二阶对称张量）、近复结构（平方为 $-1$ 倍的切丛自同构）。

\begin{exercise}
证明：$\mathrm{O}(n)/\mathrm{SO}(n)\cong\Z/2\Z$。
\end{exercise}

\begin{exercise}
证明：$\pi_1(\mathrm{SO}(n))\cong\Z/2\Z$。
\end{exercise}

\begin{exercise}
用 $\mathrm{Spin}(n)$ 表示 $\mathrm{SO}(n)$ 的万有覆盖。叙述流形的 $\mathrm{Spin}(n)$ 结构是什么。
\end{exercise}

与 G 结构有关的重要概念是可积性（integrability）。微分流形是通过把 Euclid 空间的开集用微分同胚粘贴而得到的，这一粘贴映射的微分给出切空间之间的线性映射（$\mathrm{GL}(n,\R)$ 的元素）。当粘贴映射的微分总是属于 $G$ 时，称该流形具有可积的 G 结构。换一种说法：$\R^n$ 具有标准的 G 结构；具有 G 结构的流形，在各点处的切空间（或无穷小邻域）与 $\R^n$ 连同 G 结构一起被等同；而具有可积 G 结构的流形，在各点处有限大小的邻域与 $\R^n$ 的开集连同 G 结构一起被等同。

Riemann 结构可积的充要条件是 Riemann 曲率张量为零。可积的 Riemann 流形由 $\R^n$ 的开集用等距映射粘贴而得到。这样的流形称为平坦流形；由 Bieberbach 定理可知，若它是紧的，则只能是由有限群作用取商的环面。

令
\begin{equation}
    J_{2n}:=\begin{pmatrix} 0 & -I_n \\ I_n & 0 \end{pmatrix} \tag{5.1}
\end{equation}

\begin{exercise}
证明：
\begin{equation}
    \mathrm{GL}(n,\C)\cong\left\{G\in\mathrm{GL}(2n,\R)\mid A^{-1}J_{2n}A=J_{2n}\right\} \tag{5.2}
\end{equation}
\end{exercise}

% [Tutor Note: 原文 (5.2)(5.3) 中集合的被列举元素写作 G，但条件式里的矩阵却写作 A，同一集合描述内符号混用，疑为笔误（按通常写法条件式应为 G^{-1}J_{2n}G=J_{2n} 与 G^{T}J_{2n}G=J_{2n}）；此处照原文保留，请审校确认。]

\begin{definition}
辛群（symplectic group）由
\begin{equation}
    \mathrm{Sp}(2n,\R)\cong\left\{G\in\mathrm{GL}(2n,\R)\mid A^{T}J_{2n}A=J_{2n}\right\} \tag{5.3}
\end{equation}
定义。
\end{definition}

\begin{exercise}
证明：辛结构等价于可积的 $\mathrm{Sp}(2n,\R)$ 结构。
\end{exercise}

按 Klein 的说法，辛几何是研究辛流形在辛同胚下不变的性质的学问。按 Cartan 的说法，辛几何是对可积的 $\mathrm{Sp}(2n,\R)$ 结构的研究。

\begin{exercise}
证明：
\begin{equation}
    \mathrm{U}(n)=\mathrm{O}(2n)\cap\mathrm{GL}(n,\C)\cap\mathrm{Sp}(2n,\R). \tag{5.4}
\end{equation}
而且这三个群中任意两个的交都自动包含在剩下的一个之中。
\end{exercise}

\begin{definition}
$\mathrm{U}(n)$ 结构称为近 Hermite 结构（almost Hermitian structure）。
\end{definition}

\begin{definition}
复结构可积的近 Hermite 结构称为 Hermite 结构。
\end{definition}

\begin{definition}
辛结构可积的近 Hermite 结构称为近 Kähler 结构。
\end{definition}

\begin{definition}
复结构与辛结构都可积的近 Hermite 结构称为 Kähler 结构。
\end{definition}

\begin{exercise}
证明：在辛结构、复结构与 Riemann 结构这三者中，若给定其中两个，则剩下的一个自动确定。
\end{exercise}

\begin{exercise}
举出一个例子：复结构与辛结构都可积（即流形是 Kähler 流形），但 $\mathrm{U}(n)$ 结构不可积。
\end{exercise}

Kähler 流形（或容许 Kähler 结构的复流形）构成一个非常重要的类。由 Hodge 理论，对紧 Kähler 流形，$H^q(\Omega^p_X)$ 自然地与 $(p,q)$ 次调和形式的空间 $\cH^{p,q}$ 等同，Frölicher 谱序列 (1.21) 是 $E_1$ 退化的。复共轭给出同构 $\cH^{p,q}\xrightarrow{\sim}\cH^{q,p}$，特别地 $h^{p,q}(X)=h^{q,p}(X)$。因此，紧 Kähler 流形的 Hodge 菱形左右对称，奇数阶 Betti 数为偶数。由这一点与习题 1.16 可知，Hopf 曲面不容许 Kähler 结构。

\section{模空间}

给定复流形 $X$ 时，把作为其底结构的 $C^\infty$ 级流形所能容许的复结构之集合，用复流形的同构（双全纯等价）取商所得的商集合，在良好情形下具有复流形（或其推广）的结构。这称为 $X$ 的复结构的模空间。

\begin{remark}
把「按所有双全纯等价取商」换成「只按与恒等映射同痕（isotopic）的全纯等价取商」所得的空间，称为 Teichmüller 空间。Teichmüller 空间具有以映射类群为结构群的、模空间上的主丛结构。
\end{remark}

设 $V$ 为 $n$ 维复向量空间，$\Lambda$ 为秩 $2n$ 的自由 Abel 群。

\begin{exercise}
证明：对任意 $n$ 维复环面 $X$，存在单射群同态 $i\colon\Lambda\to V$，使得 $X\cong V/i(\Lambda)$。
\end{exercise}

把单射同态的集合记作 $\Hom(\Lambda,V)^\circ\subset\Hom(\Lambda,V)$。

\begin{exercise}
证明：对 $i,i'\in\Hom(\Lambda,V)^\circ$，$V/i(\Lambda)\cong V/i'(\Lambda)$ 的充要条件是：存在 $g\in\mathrm{GL}(V)$ 与 $h\in\Aut(\Lambda)$，使得 $i'=g\circ i\circ h$。
\end{exercise}

因此，实 $2n$ 维环面上可容许的复结构的等价类之集合，由
\begin{equation}
    \mathrm{GL}(V)\backslash\Hom_\Z(\Lambda,V)^\circ/\Aut(\Lambda) \tag{6.1}
\end{equation}
给出。

\begin{exercise}
证明：$\mathrm{GL}_n(\C)\backslash\Hom(\Lambda,V)^\circ$ 自然地与 $2n$ 维复向量空间中 $n$ 维复子空间的集合 $\mathrm{Gr}_n(\C^{2n})$ 的某个子集等同，并具有自然的 $n^2$ 维复流形结构。
\end{exercise}

\begin{exercise}
证明：当 $n=1$ 时，$\mathrm{GL}_n(\C)\backslash\Hom(\Lambda,V)^\circ\cong\C\setminus\R$，且 $\mathrm{GL}(V)\backslash\Hom_\Z(\Lambda,V)^\circ/\Aut(\Lambda)\cong\mathbb{H}/\mathrm{SL}_2(\Z)$。这里 $\mathbb{H}:=\{z\in\C\mid\operatorname{Im}z>0\}$ 是上半平面。
\end{exercise}

\begin{exercise}
证明：对 $n\ge2$，集合 $\mathrm{GL}(V)\backslash\Hom_\Z(\Lambda,V)^\circ/\Aut(\Lambda)\cong\mathrm{Gr}_n(\C^{2n})/\mathrm{GL}_{2n}(\Z)$ 关于商拓扑不是 $T_1$ 的（从而也不是 Hausdorff 的）。
\end{exercise}

另一方面，已知 $n$ 维极化 Abel 簇的同构类之集合自然具有 $\frac{1}{2}n(n+1)$ 维复解析空间的结构。对 $n$ 维环面上的任意线丛（line bundle）$L$，请注意 $L$ 的变形的障碍类所在空间的维数为 $\dim\Ext^2(L,L)=\binom{n}{2}=n^2-\frac{1}{2}n(n+1)$。

\begin{remark}
辛结构的模空间也可以同样地定义，但照原样并不是太好的概念。在镜像对称的语境中，人们认为辛结构的模空间会受到所谓量子修正（quantum correction），而修正后的模空间的令人满意的数学表述尚未给出。
\end{remark}

\section{一切都是 Lagrange 子流形}

相空间是给位形空间添上动量而「加倍」得到的；其中动量为零的子空间 $L$ 可自然地与位形空间等同。换言之，$M=T^*N$ 的零截面 $L$ 可与 $N$ 等同。这个 $L$ 具有一个显著性质：$M$ 的辛形式 $\omega$ 限制到 $L$ 上为零，即 $\omega|_L=0$。而且，在具有这种性质的 $M$ 的子流形中，$L$ 关于包含关系是极大的。具有这两条性质的 $M$ 的子流形称为 Lagrange 子流形：

\begin{definition}
辛流形 $(M,\omega)$ 的子流形 $L$ 称为 Lagrange 子流形（Lagrangian submanifold），是指 $\omega|_L=0$ 且 $\dim L=\frac{1}{2}\dim M$ 成立。
\end{definition}

「一切都是 Lagrange 子流形」这一主张，称为 Weinstein 的辛教条（Weinstein's symplectic creed）[Wei81]。纵然「一切」这种说法略有夸张，辛几何中的许多重要问题都能用 Lagrange 子流形的语言写出，这一点是确实的。\footnote{在移位辛几何（shifted symplectic geometry）中，移位辛结构（shifted symplectic structure）是移位 Lagrange 结构（shifted Lagrangian structure）的特例（例如参见 [Toë14, Remarks 5.3]）。}

例如，辛同胚
\begin{equation}
    f\colon(M_1,\omega_1)\to(M_2,\omega_2) \tag{7.1}
\end{equation}
可通过考虑其图像，看作直积 $(M_1\times M_2,\pi_1^*\omega_1-\pi_2^*\omega_2)$ 的 Lagrange 子流形的特例。特别地，当辛流形 $M_1$ 与 $M_2$ 都是另一个流形 $N$ 的余切丛、且该 Lagrange 子流形作为直积 $N\times N$ 上的函数 $F$ 的微分 $dF$ 的图像给出时，称 $F$ 为该变换的生成函数。具体地说，设 $M_1$ 的坐标为 $(q,p)$、$M_2$ 的坐标为 $(Q,P)$，则用生成函数 $F(q,Q)$，变换前后变量的关系可写作
\begin{align}
    p_i &= \frac{\partial F}{\partial q_i}, \tag{7.2} \\
    P_i &= -\frac{\partial F}{\partial Q_i} \tag{7.3}
\end{align}
由于 $F$ 是 $q$ 与 $Q$ 的函数，先解 (7.2) 把 $Q$ 表为 $q$ 与 $p$ 的函数，再由 (7.3) 确定 $P$ 为 $q$ 与 $p$ 的函数。

\begin{exercise}
用 Lagrange 子流形来记述矩映射。\footnote{对这一问题的「良好」解答，可以借助移位辛结构的语言给出。矩映射 $\mu\colon X\to\fg^*$ 所诱导的商栈的态射 $\mu\colon[X/G]\to[\fg^*/G]\simeq T^*BG[1]$，关于 $T^*BG[1]$ 的标准的 1-移位辛结构具有典范的 Lagrange 结构。同样地，设 $\lambda\in\fg^*$ 的稳定化子为 $G_\lambda$，则自然嵌入 $BG_\lambda\hookrightarrow[\fg^*/G]$ 具有典范的 Lagrange 结构。在这些纤维积 $[X/G]\times_{[\fg^*/G]}BG_\lambda$ 上赋予典范的 0-移位辛结构，所得之物就是 $X$ 在 $\lambda$ 处的（导出）辛商。}
\end{exercise}

% CHUNK c03 DONE pages p009–p013
